\begin{proof}
Fix a displayed trivializing open \(U\), restrict \(h\) to its inverse
image, and transport it through the analytic isomorphism with
\(U\times_KF\).  Ordinary-open restriction and transport through an
analytic isomorphism preserve the zero-current equality.  Choose
\(p\in D(K)\).  If
\(x\in U(K)\), then \(\mathscr H(x)=K\).  The exact fibre-slicing theorem
\Cref{prop:p062-strict-product-fibre-slicing-v4-research} therefore makes
the restriction \(h_x\) to \(F\) a continuous zero-current function over
\(K\).  By \Cref{cor:ch19-continuous-zero-current-harmonic}, \(h_x\) is
classically pluriharmonic.  It attains a maximum because \(F\) is proper.
The maximum locus is locally open by
\Cref{prop:classical-psh-calculus}, and it is nonempty and closed; since
\(F\) is connected, \(h_x\) is constant.

Let \(s_{U,p}:U\to U\times_KF\) be the constant section and put
\(b_U:=h\circ s_{U,p}\).  The preceding constancy gives
\[
                    h(y)=\operatorname{pr}_U^*b_U(y)
                    \tag{\(\spadesuit\)}
\]
at every rigid point \(y\) of \(U\times_KF\): its image in \(U\) is a
rigid point, hence a \(K\)-rational point because \(K\) is algebraically
closed.  Rigid points are dense locally in a topologically separated strict
analytic space \cite[Proposition~2.1.15]{Berk12}.  Applying this to the
good strict product and using continuity extends \((\spadesuit)\) to all
of \(U\times_KF\).  This use of the section identifies functions only; it
does not pull a current back along the section.

Choose a compactly supported ordinary CLD-smooth \((1,1)\)-form \(\eta\)
on \(F\) with \(\int_F\eta=1\).  One obtains it by taking a nonzero
nonnegative top form in a strict coordinate disc about \(p\), supported in
a relatively compact strict annulus, and normalizing; compare the identical
construction in \Cref{prop:p062-open-disc-descent} and
\cite[Proposition~11.1 and Example~11.7]{GJR26}.  For every compactly
supported ordinary smooth \((m-1,m-1)\)-form \(\beta\) on \(U\), the
top fibre degree gives
\[
 \mathrm d'\mathrm d''\bigl(\operatorname{pr}_U^*\beta
       \wedge\operatorname{pr}_F^*\eta\bigr)
   =\operatorname{pr}_U^*(\mathrm d'\mathrm d''\beta)
       \wedge\operatorname{pr}_F^*\eta.
\]
Pairing this with the zero current of \(h\), using \((\spadesuit)\), and
integrating over the fibre gives
\[
 0=\int_{U\times_KF}h\,
       \operatorname{pr}_U^*(\mathrm d'\mathrm d''\beta)
       \wedge\operatorname{pr}_F^*\eta
   =\int_U b_U\,\mathrm d'\mathrm d''\beta.
\]
The continuous-multiplier fibre-integration identity is the compact-support
approximation and projection-formula argument in
\Cref{prop:p062-open-disc-descent}.  Hence
\(\mathrm d'\mathrm d''[b_U]=0\).  If \(\mathrm d'[h]=0\), the same
calculation with a compactly supported ordinary smooth \((m-1,m)\)-form
on \(U\) gives \(\mathrm d'[b_U]=0\).

On overlaps, the local pullback identities and the nonempty fibres force
the functions \(b_U\) to agree.  They glue to a unique continuous function
\(b\) on \(B\).  Their zero-current identities, and likewise their
first-current identities when present, glue by
\Cref{prop:ch14-currents-sheaf}; the stated hypotheses on \(B\) are those
of that proposition.  This proves all assertions, including uniqueness.
\end{proof}
